\section{Plan de mejora reorganizado por prioridad y dependencias}
\label{sec:plan}
\textbf{Fecha: 18 de octubre de 2026. Estado: bloques I--V con implementación local; validaciones humanas, musicales y de rendimiento con modelos reales pendientes.}
Esta secuencia sustituye el orden de ejecución anterior. Los cambios 1--8 ya
realizados conservan su registro; los puntos originales 9--14 y los IDs
E01--E12 se mantienen como referencias de trazabilidad. No son listas de
trabajo adicionales: cada pendiente se integra en una o más actividades P.

Los registros de ejecución están en las secciones~\ref{sec:bloque1},
\ref{sec:bloque2}, \ref{sec:bloque3}, \ref{sec:bloque4} y~\ref{sec:bloque5}.
P01--P03 y P05--P08 cuentan con verificación local; P04 tiene protocolo y línea
base simulada. Las anotaciones y medidas reales permanecen pendientes.
P09/P11 están implementados; P10 tiene contrato y pruebas, con calibración
humana pendiente; P12 tiene comparación y política, con utilidad real pendiente.
P13--P15 están implementados; P16 tiene instrumentación y perfil local con
índices reales/modelos de control, pendiente de carga y memoria con modelos
reales. P17--P18 tienen implementación y pruebas locales; P19 tiene adaptador
opcional y P20 regresiones integradas, con integración real pendiente.
E13 se añadió tras el bloque II y se
integra en actividades existentes, sin crear un sexto bloque.
La revisión posterior confirmó E14 (versionado por reordenación), E15
(validación JSON omitida en aclaraciones) y E16 (fallos de registro en CP1252).
Se reasignaron y corrigieron en la continuación: E14 al bloque II/P07
(identidad del plan), E15 al bloque II/P05/P08 (estado y planificación) y
E16 al bloque I/P01/P02 (portabilidad e inicialización). No crean un bloque
adicional ni desplazan P09--P12. E14 también tiene regresiones en P11/P12.
La evidencia histórica está en \texttt{validation/revision\_block2/}; las
correcciones y nuevas pruebas, en \texttt{validation/block3/}.

\paragraph{Revisión posterior al bloque V.}
La auditoría confirmó E13 y aprobaciones \texttt{legacy} sin citas. La corrección
posterior verifica correspondencia literal o equivalencias revisadas, exige el
alcance completo de negaciones compuestas y rechaza traducciones no verificadas.
El juez exige JSON con citas y se abstiene ante formatos antiguos persistentes;
el orquestador tampoco selecciona una aprobación \texttt{legacy}.
Estas correcciones locales quedan en bloque II/P07--P08 y bloque III/P10;
no demuestran fidelidad de lenguaje arbitrario ni calidad musical.
Ahora la prioridad es anotar y medir P04/P10/P12 y rendimiento real P16.
La captura previa está en \code{validation/revision_previous_blocks/audit.json};
la verificación posterior, en \code{validation/pending_fixes/} y
sección~\ref{sec:correcciones_pendientes}. E14--E16 siguen corregidos;
no se reabren sin una regresión demostrada. E17--E19 permanecen en sus actividades
de ingesta, retención y modalidades indicadas en el bloque V.
La continuación autónoma implementa E17 (cuotas, retención/cuarentena, backup y
restauración) y E19 (mapas independientes e instrumentales declarados). E18 tiene
entorno opcional fijado e ingesta técnica real; se corrigió una incompatibilidad
Nomic/Transformers 5 y se fijan revisiones de pesos/código. La sección~\ref{sec:autonomo}
registra pruebas reales de workers y la revisión visual/PDF. Continúan pendientes
el proveedor Groq, las anotaciones humanas y la calidad de transcripción musical.

Se establecen cinco bloques de cuatro actividades. Primero se hace reproducible
la ejecución y se registra una línea base; después se corrige la interpretación
y el control; luego se revisan fallos y refinamiento; finalmente se amplían
producto, operación e ingesta. No se añaden cuatro funciones nuevas sobre una
base cuyo comportamiento todavía no pueda comprobarse.

\subsection{Bloque I: ejecución verificable y línea base}
\textbf{Prioridad crítica. Estado: P01--P03 implementados; P04 disponible en modo control, con medición real pendiente.}
Estas actividades permiten reproducir los problemas y comparar los bloques
posteriores contra un estado conocido.

\begin{description}
\item[P01. Entorno limpio y dependencias reproducibles.]
Separar dependencias directas y de pruebas, resolver versiones compatibles,
crear configuración de ejemplo e instrucciones de instalación y comprobar
imports en un entorno nuevo, sin reutilizar la carpeta temporal de esta sesión.
\textbf{Aceptación:} instalación e imports reproducibles con versiones
inequívocas. \textbf{Origen:} punto 12 y E11. \textbf{Dependencias:} ninguna.

\item[P02. Contrato HTTP, arranque y configuración.]
Validar cuerpo JSON y tipos antes de operar sobre ellos; definir estados HTTP,
inicialización web/WSGI y rutas estables. Separar configuración de desarrollo
y ejecución del servicio. Probar con Flask real.
\textbf{Aceptación:} entradas inválidas tienen respuestas controladas y el motor
se inicializa en los puntos de entrada soportados.
\textbf{Origen:} punto 12 y E11. \textbf{Dependencia:} P01.

\item[P03. Validación numérica y coherencia de datos.]
Comprobar dimensiones, normas, valores finitos, vectores nulos, IDs y
correspondencia entre índices, mapa y catálogo al cargar y antes de buscar.
\textbf{Aceptación:} los datos incompatibles y consultas degeneradas no producen
rankings aparentemente válidos. \textbf{Origen:} parte numérica de E10.
\textbf{Dependencia:} P01.

\item[P04. Línea base y protocolo de evaluación.]
Fijar corpus y consultas líricas, acústicas, mixtas, literales y negativas,
etiquetas de referencia, registros y mediciones. Capturar el comportamiento
actual antes de las siguientes modificaciones; separar modos simulado y real.
\textbf{Aceptación:} evaluación repetible que registra relevancia, cumplimiento,
estabilidad, latencia, llamadas y errores; declara métricas aún no ejecutadas
si faltan modelos o acceso al proveedor.
\textbf{Origen:} punto 13 y base experimental de E07.
\textbf{Dependencias:} P01--P03 para la línea base integrada.
\end{description}

\subsection{Bloque II: intención, modalidades y estado coherente}
\textbf{Prioridad alta. Estado: implementación y pruebas locales completas.} Se define qué se intenta recuperar, qué evidencia lo
respalda y quién modifica las decisiones. Es la base de una degradación segura.

\begin{description}
\item[P05. Propiedad única del estado y las exclusiones.]
El juez devuelve una decisión y solo el orquestador la aplica. Separar estado
serializable y recursos efímeros; vincular decisiones a candidato y versión
de plan, conservando el aislamiento entre solicitudes.
\textbf{Aceptación:} un único punto de escritura de exclusiones y pruebas de
invariantes y concurrencia. \textbf{Origen:} E08.
\textbf{Dependencias:} P02 y P04.

\item[P06. Rutas líricas, acústicas e instrumentales.]
Definir evidencia y estado de evaluación por modalidad. Permitir búsquedas
solo audio y canciones instrumentales sin exigir letra. No dirigir cambios
acústicos a partir de un rechazo lírico sin diagnóstico justificado.
\textbf{Aceptación:} un instrumental no se excluye por carecer de letra y una
aprobación lírica no se presenta como validación del sonido.
\textbf{Origen:} E01 y alcance de validación del punto 11.
\textbf{Dependencias:} P03--P05.

\item[P07. Restricciones explícitas e hipótesis del modelo.]
Registrar la procedencia de positivos y negativos; permitir reparar una
modalidad omitida o un negativo inventado sin borrar restricciones del usuario.
Definir cuándo un rechazo depende del plan y puede reconsiderarse.
\textbf{Aceptación:} corregir un error del router no relaja silenciosamente
una restricción explícita. \textbf{Origen:} E03 y revisión de exclusiones de E01.
\textbf{Dependencias:} P05--P06.

\item[P08. Fallback fiel a la intención.]
Usar el contrato y la procedencia de restricciones para preservar lo conocido
cuando falla el router; tratar explícitamente idioma y modalidades inciertas.
No asumir que copiar la consulta en las dos ramas conserva su significado.
\textbf{Aceptación:} fallos del router en consultas con negaciones o en español
no eliminan silenciosamente restricciones; la degradación queda declarada.
\textbf{Origen:} E02. \textbf{Dependencias:} P06--P07.
\end{description}

\subsection{Bloque III: fallos, evaluación y refinamiento fundamentado}
\textbf{Prioridad alta. Implementación local realizada; validación real pendiente en P10/P12.} Este bloque corrige regresiones de disponibilidad y
evalúa si conviene seguir refinando, con la línea base ya preparada.

\begin{description}
\item[P09. Recuperación parcial y continuidad ante fallos.]
Definir degradación por modalidad, evidencia del plan usado y continuidad con
candidatos existentes si falla el refinador. Ajustar pesos solo bajo una política
explícita y probada.
\textbf{Aceptación:} una rama fallida no borra sin explicación resultados útiles
de la otra, ni su ausencia se convierte en aprobación.
\textbf{Origen:} E05. \textbf{Dependencias:} P05--P08.
\textbf{Estado:} implementado, con degradación por rama y restauración del
plan/ranking anterior si falla toda la recuperación reformulada.

\item[P10. Juez y evidencia lírica suficiente.]
Evaluar salida estructurada, abstención, reintento limitado y selección de
fragmentos más allá de los primeros 3000 caracteres. Calibrar contra juicios
humanos y mantener estados por candidato.
\textbf{Aceptación:} medir falsos rechazos/aprobaciones y cubrir evidencia al
final de una letra sin marcar todo el ranking como evaluado.
\textbf{Origen:} E06 y parte de evidencia del punto 11.
\textbf{Dependencias:} P04--P07.
\textbf{Estado:} contrato estructurado, citas, abstención, cobertura y reintento
limitado implementados; tasas de error contra juicios humanos aún no medidas.

\item[P11. Presupuestos de recuperación y evaluación.]
Separar descartes por metadatos, llamadas al juez, refinamientos y ampliación
del ranking. Determinar qué significa agotamiento sin confundirlo con ausencia
de canciones válidas en el catálogo completo.
\textbf{Aceptación:} metadatos incompletos no consumen todos los intentos útiles
y las búsquedas adicionales tienen un límite verificable.
\textbf{Origen:} parte de presupuesto de E10.
\textbf{Dependencias:} P03, P06 y P09--P10.
\textbf{Estado:} implementado con presupuestos separados y ampliación limitada;
se informa agotamiento de ventana, sin afirmar agotamiento del catálogo.

\item[P12. Decidir cuándo el refinamiento aporta valor.]
Comparar baseline corregido y flujo con refinamiento sobre las mismas consultas,
con coste y calidad. Habilitarlo, limitarlo o retirarlo según evidencia y tipo
de consulta, manteniendo explícito que el juez actual no escucha audio.
\textbf{Aceptación:} resultados comparables con dispersión y criterio de
activación; no declarar superioridad si no aparece en la evaluación.
\textbf{Origen:} E07 y punto 13.
\textbf{Dependencias:} P04 y P09--P11.
\textbf{Estado:} comparación pareada y política experimental implementadas;
refinamiento desactivado por defecto. La simulación no demuestra relevancia;
decisión de activación rutinaria pendiente de evidencia real y humana.
\end{description}

\subsection{Bloque IV: resultados útiles y rendimiento medido}
\textbf{Prioridad media. P13--P15 implementados; P16 con perfil local y verificación real pendiente.}
Mantener P10/P12 como verificaciones abiertas, sin declararlas cerradas por
avanzar en capacidades independientes. La funcionalidad se amplía después de estabilizar la
semántica y el manejo de errores.

\begin{description}
\item[P13. Puntuación interpretable y pesos configurables.]
Exponer contribuciones por modalidad, revisar la fusión y sustituir el porcentaje
por un score identificado correctamente. Permitir ajustar pesos con validación.
\textbf{Aceptación:} la UI y API explican el score sin presentarlo como
probabilidad y los pesos válidos producen resultados reproducibles.
\textbf{Origen:} punto 9. \textbf{Dependencias:} P04, P06 y P09--P12.
\textbf{Estado:} fusión por rangos ponderados, contribuciones, pesos validados
y API v2 sin porcentaje implementados. No demuestra mayor relevancia musical.

\item[P14. Filtros exactos de metadatos.]
Añadir artista, género y rango de años con semántica explícita antes de la
selección final; distinguir filtro duro de preferencia semántica.
\textbf{Aceptación:} los resultados cumplen los filtros y las combinaciones
sin coincidencias se informan correctamente.
\textbf{Origen:} punto 10. \textbf{Dependencias:} P03 y P07--P11.
\textbf{Estado:} filtros exactos antes de fusión/juez, rango inclusivo y opciones
de catálogo implementados; selección completa para el catálogo actual.

\item[P15. Evidencia y validación visibles por resultado.]
Mostrar fragmentos relevantes de letras, contribuciones, modalidades evaluadas,
advertencias y alternativas sin atribuir aprobación a canciones no juzgadas.
\textbf{Aceptación:} cada tarjeta refleja fielmente la evidencia disponible y
el alcance de su evaluación. \textbf{Origen:} punto 11, complementando P06/P10.
\textbf{Dependencias:} P10 y P13--P14.
\textbf{Estado:} fragmentos con procedencia, contribuciones y alcance por
candidato visibles; navegador verificado con fixture de control, sin validar sonido.

\item[P16. Optimización con mediciones de carga real.]
Medir compilación del grafo, caches, codificación, memoria y contención.
Valorar grafo reutilizable, batch y tamaño de caches solo donde haya beneficio;
comparar con una máquina de estados simple si procede.
\textbf{Aceptación:} publicar latencia fría/caliente, consumo y carga concurrente
sin perder invariantes funcionales.
\textbf{Origen:} rendimiento de E04/E09 y E07.
\textbf{Dependencias:} P04, P12 y estado funcional del bloque IV.
\textbf{Estado:} tiempos, contención y cachés instrumentados; perfil de índices
reales con encoders/juez de control completado. Inferencia, RSS/GPU y optimización
real siguen pendientes. Se conservan grafos por solicitud sin nuevas cachés globales.
\end{description}

\subsection{Bloque V: catálogo actualizable y cierre integral}
\textbf{Implementación local realizada; validación externa pendiente.} La ingesta se aborda al contar
con un catálogo coherente y una operación verificable.

\begin{description}
\item[P17. Snapshots del catálogo e invalidación de caches.]
Versionar catálogo, mapa e índices, publicar actualizaciones coordinadas y
definir invalidación o reinicio controlado. Evitar mutaciones de recursos
compartidos por consumidores.
\textbf{Aceptación:} ninguna consulta mezcla artefactos de versiones distintas
y una actualización fallida conserva un snapshot utilizable.
\textbf{Origen:} coherencia de E09. \textbf{Dependencias:} P03, P05 y P16.
\textbf{Estado:} snapshots inmutables, publicación transaccional e invalidación por
versión implementados y probados; véase sección~\ref{sec:bloque5}.

\item[P18. Decidir sobre checkpoints y progreso en vivo.]
Revisar qué capacidades de C se necesitan realmente. Si se requieren, añadir
persistencia y streaming sobre estado serializable; si no, registrar por qué
se difieren. No reinstaurarlas solo por haber existido en un experimento.
\textbf{Aceptación:} decisión documentada; si se implementan, reanudar o
transmitir progreso no mezcla sesiones ni expone estados incoherentes.
\textbf{Origen:} funciones de E04 y serialización de E08.
\textbf{Dependencias:} P05, P09 y P16.
\textbf{Estado:} decisión documentada: búsquedas sin checkpoints; ingesta con cola
SQLite, propietario, lease y progreso consultable. No se reanuda un trabajo vencido.

\item[P19. Ingesta de canciones en segundo plano.]
Incorporar YouTube, FFmpeg, subtítulos y Whisper cuando corresponda, cálculo de
embeddings y seguimiento de trabajos; publicar canciones mediante el mecanismo
de catálogo de P17.
\textbf{Aceptación:} añadir una canción no bloquea la web, informa progreso/error
y solo publica datos y vectores completos y coherentes.
\textbf{Origen:} punto 14. \textbf{Dependencias:} P01--P03 y P17;
P18 solo si se elige reutilizar su mecanismo de progreso.
\textbf{Estado:} adaptador y cola implementados, probados con dependencias de
control. Pendientes ejecución real, entorno opcional fijado e instrumentales sin
transcripción (E18--E19). No se declara cerrada la aceptación integral.

\item[P20. Validación integrada y cierre de pendientes.]
Ejecutar recorridos completos de web y consola, ingesta si está habilitada,
concurrencia, restricciones, fallos y calidad con el entorno reproducible.
Actualizar instalación, documentación y decisiones de activar o diferir funciones.
\textbf{Aceptación:} informe reproducible que separa pruebas simuladas y reales,
beneficios medidos, limitaciones y pendientes no cerrados.
\textbf{Origen:} cierre integrado de puntos 12--14 y E07/E11.
\textbf{Dependencias:} actividades previas implementadas o diferidas explícitamente.
\textbf{Estado:} regresiones locales integradas; calidad musical, inferencia real,
recuperación ante apagado y operación prolongada siguen pendientes.
\end{description}

\subsection{Correspondencia con el listado anterior}
\small
\begin{longtable}{p{0.24\textwidth}p{0.65\textwidth}}
\toprule Referencia previa & Nueva actividad de ejecución\\\midrule\endhead
Punto 9 & P13.\\
Punto 10 & P14.\\
Punto 11 & P06, P10 y P15.\\
Punto 12 & P01--P02 y P20.\\
Punto 13 & P04, P12 y P20.\\
Punto 14 & P19, con P17 como requisito de publicación.\\
E01 & P06--P07.\\
E02 & P08.\\
E03 & P07.\\
E04 & P16 y P18.\\
E05 & P09.\\
E06 & P10 y P15.\\
E07 & P04, P12, P16 y P20.\\
E08 & P05 y P18.\\
E09 & P16--P17.\\
E10 & P03 y P11.\\
E11 & P01--P02 y P20.\\
E12 & Seguimiento transversal; no añade una actividad duplicada.\\
E13 & P10/P12; P15 si se necesita corrección visible de restricciones.\\
E14 & Bloque II/P07; regresiones en P11/P12. Corregido en la continuación.\\
E15 & Bloque II/P05/P08. Corregido en la continuación.\\
E16 & Bloque I/P01/P02. Corregido y probado en CP1252 y UTF-8.\\
\bottomrule
\end{longtable}
\normalsize

\textbf{Regla de avance.} Cada actividad registra implementación, evidencia y
limitaciones antes de cerrarse. El orden puede ajustarse ante un defecto
bloqueante demostrado; cualquier cambio conserva IDs y justifica su dependencia.
Las pruebas y la documentación se actualizan en todos los bloques, no se aplazan
hasta P20. Las medidas de relevancia real necesitan modelos/proveedor y referencia
humana; si faltan, esa parte permanece pendiente y no se sustituye por un éxito
de pruebas simuladas.
